\documentclass{resume}
\usepackage[left=0.5in,top=0.6in,right=0.5in,bottom=0.6in]{geometry} % Document margins

\name{Jincheng (Eric) Huang}
\address{700 19th St NW, Washington DC, 20431}
\address{ehuang@imf.org}

\begin{document}

\begin{rSection}{Education}
{\bf University of Pennsylvania} \hfill {\em Sept 2017 - May 2023 (expected)}
\\M.S. \& Ph.D. in Economics 

{\bf University of Wisconsin-Madison} \hfill {\em Sept 2011 - May 2015} 
\\B.A. in Economics (with Honors), Mathematics and Computer Science

\end{rSection}

\begin{rSection}{Research}
\textbf{Precautionary Mismatch} with Xincheng Qiu (Job Market Paper)  \\
\begin{small}
How does wealth affect the allocative efficiency of talents? We study this question using a model with worker and firm heterogeneity, search frictions and incomplete markets. In the model, workers and firms jointly face a trade-off between the speed of match formation and the profit from the match, and are thus willing to accept a certain degree of mismatch. For workers having little wealth while searching for jobs, this trade-off is weighed in favor of speed due to precautionary motive, leading to a higher degree of skill mismatch. We call this phenomenon ``precautionary mismatch''. We estimate that the output loss due to skill mismatch in the current U.S. working population is roughly 3\%. Without precautionary motive, aggregate labor productivity would increase by 1.6\%, mainly due to improvement in the allocation of high-skilled liquidity-constrained workers. 
\end{small} 
\vspace{5pt}
\\
\textbf{Saving for Unexpected Expenses} (work in progress)
\end{rSection}

\begin{rSection}{Work Experience}
\begin{rSubsection}{Federal Reserve Bank of St. Louis}{Aug 2021}{Dissertation Intern}{}
    \item Presented dissertation research project at the Bank's workshops and seminars.
    \item Held follow-up feedback sessions with the Bank's economists regarding research project.
\end{rSubsection}

\begin{rSubsection}{The Wharton School}{Jan 2020 - Sept 2021}{Graduate Research Assistant}{}
    \item Cleaned and analyzed large data sets on housing transactions, land parcels and neighborhood features.
    \item Developed causal inference framework to study the environmental and economic impacts of highway pollution.
    \item Arranged weekly feedback sessions with principal investigators.
\end{rSubsection}

\begin{rSubsection}{Federal Reserve Bank of New York}{July 2015 - July 2017}{Research Analyst \& Senior Research Analyst}{}
    \item Assisted the Bank's economists with implementation and examination of research designs.
    \item Provided statistical analysis and research support for Fed policy briefings and balance sheet simulations.
    \item Investigated and reconciled stress test models to predict mortgage default and prepayment risks.
\end{rSubsection}

\end{rSection}

\begin{rSection}{Teaching Experience}
Econ 10: \textit{Introductory Economics for Business Students}, University of Pennsylvania
\\ Teaching assistant for Gizem Saka \hfill {\em Fall 2018, 2019}

\\Econ 2: \textit{Introductory Macroeconomics}, University of Pennsylvania
\\ Teaching assistant for Luca Bossi \hfill {\em Spring 2019}
\end{rSection}

\pagebreak

\begin{rSection}{Honors \& Scholarships}
Graduate Research Fellowship, National Science Foundation \hfill {\em 2017 - Present} \vspace{1pt}\\
Xingmei Zhang \& Yongge Dai Fellowship, University of Pennsylvania \hfill {\em 2018 - 2019} \vspace{1pt}\\
Vault Award, Federal Reserve Bank of New York \hfill {\em Oct 2016} \vspace{1pt}\\
Meek Bishop Award in Economics, University of Wisconsin-Madison \hfill {\em Dec 2014} \vspace{1pt}\\
L. G. Ammerman Award in Economics, University of Wisconsin-Madison \hfill {\em May 2014} \vspace{0pt}
\end{rSection}

\begin{rSection}{Conference \& Seminar Presentations}
2022: Penn, WashU EGSC \vspace{1pt}\\
2021: Penn, ESPE 2021 Annual Conference, CES 2021 Conference, Warwick Economics PhD Conference, SED Annual Meeting, St. Louis Fed
\end{rSection}

\begin{rSection}{Professional Activities}
Referee Services: Macroeconomic Dynamics
\end{rSection}

\begin{rSection}{Technical Skills}
Stata, MATLAB, \LaTeX, Python, Julia, Java, Git
\end{rSection}

\end{document}
